\documentclass[12pt]{article}
\usepackage[utf8]{inputenc}
\usepackage[T1]{fontenc}
\usepackage{libertinus}
\usepackage{libertinust1math}
\usepackage[protrusion=true,expansion=true]{microtype}
\usepackage[margin=1.4in]{geometry}
\usepackage{setspace}
\usepackage{booktabs}
\usepackage[colorlinks=true,linkcolor=black,urlcolor=black]{hyperref}
\setlength{\parindent}{0pt}
\setlength{\parskip}{0.6em}

\title{Contents of \textit{Samlede Skrifter}, Trettende Bind\\[0.4em]
  \large Georg Brandes (Copenhagen: Gyldendal, 1903)}
\author{}
\date{}

\begin{document}
\maketitle

\noindent
The thirteenth volume of Brandes's collected works is a supplementary
volume (\textit{Supplementbind}) gathering his principal contributions
to the \textit{Tro og Viden} controversy of the 1860s, together with
several longer critical essays. The philosophical core is the two
sections on the faith--knowledge dispute: \textit{Striden om Tro og
Viden} (pp.~5--42) and \textit{Dualismen i vor nyeste Filosofi}
(pp.~43--84). A full OCR text of the volume is available at
Project Runeberg (\url{https://runeberg.org/gbsamskr/13/}).

\bigskip
\begin{tabular}{p{9cm}r}
\toprule
\textbf{Title} & \textbf{Page} \\
\midrule
Indledning & 1 \\
\midrule
\textit{Striden om Tro og Viden} & 5 \\
\quad Samvittighed og Videnskab & 7 \\
\quad Svar til Hr.\ Schmidt & 11 \\
\quad Filosofi og Teologi & 13 \\
\quad En sidste Replik & 17 \\
\quad Om Selvansvarlighed & 22 \\
\quad Nationalitetsideens Gyldighed & 26 \\
\quad Indledning til den rationelle Etik & 28 \\
\quad Gensvar til Dr.\ Heegaard & 34 \\
\quad Grundideernes Logik & 40 \\
\midrule
\textit{Dualismen i vor nyeste Filosofi} & 43 \\
\quad Indledning & 43 \\
\quad Den absolute Uensartethed & 54 \\
\quad Det Antirationelle & 67 \\
\quad Dualismens F\o{}lger & 71 \\
\midrule
Om den gode Vilje som Magt i Videnskaben & 85 \\
Nielsen og Grundtvig & 93 \\
Et Svar & 98 \\
F\ae{}drelandets Hr.\ 99 & 102 \\
R.\ Nielsens L\ae{}re om Tro og Viden & 106 \\
Biskop Jens Paludan-M\"uller & 110 \\
Filosofiens historiske Udvikling & 115 \\
Den lille R\o{}dh\ae{}tte & 124 \\
Breve til en Landsbypr\ae{}st & 126 \\
\midrule
Den franske \AE{}stetik i vore Dage & 137 \\
Den herskende Evne & [?] \\
Kritiken & [?] \\
\bottomrule
\end{tabular}

\bigskip
\section*{Notes on individual pieces}

\textbf{Indledning} (pp.~1--4). Brandes's preface to the volume,
dated April 1903, explaining the rationale for collecting these
earlier polemical writings.

\textbf{Striden om Tro og Viden} (pp.~5--42). A sequence of
polemical articles from the controversy of the mid-1860s, most
originally published in newspapers and journals. The first piece,
\textit{Samvittighed og Videnskab} (pp.~7--10), is a review of
A.-C.\ Larsen's book of the same title (1866), a theological
response to Nielsen. Brandes notes that Nielsen's standpoint
invites attack from two directions: from theologians he is too
radical, and from secularists he remains too theological.
\textit{Gensvar til Dr.\ Heegaard} (pp.~34--39) is a direct
response to a Nielsen disciple.

\textbf{Dualismen i vor nyeste Filosofi} (pp.~43--84). The
centerpiece of the volume: Brandes's sustained philosophical
attack on Nielsen's absolute separation of faith and knowledge.
Structured in four parts. Corresponds to the 1866 monograph
\textit{Dualismen i vor nyeste Philosophie} (note the modernized
spelling in the collected works). The most philosophically
substantial piece in the volume.

\textbf{R.\ Nielsens L\ae{}re om Tro og Viden} (pp.~106--109).
A compact later statement directly on Nielsen's doctrine.

\textbf{Den franske \AE{}stetik i vore Dage} (pp.~137ff.). Longer
critical essay on French aesthetics; less directly relevant to
the faith--knowledge debate.

\bigskip
\textit{Note}: page numbers for the final three entries
(\textit{Den herskende Evne}, \textit{Kritiken}) are not yet
confirmed from the OCR source and are marked~[?].

\end{document}
